%%%PAGE 7 rot=0 legibility=good summary=电子轨道跃迁（一群/一个、光照/轰击跃迁与电离）；天然放射现象α β γ对比表；核反应三变四恒与质量；α/β/γ衰变；人工转变（卢瑟福发现质子、查德威克发现中子）
%%%BLOCK id=007-01 ch=17 sec=氢原子能级 kind=knowledge alt=none cont=none y=0.01-0.23
\textbf{电子轨道跃迁}

一群\,/\,一个\quad $n$能级氢原子向低能级跃迁时，辐射光子的 能量/频率/波长/光谱 有\quad $\left\{\begin{array}{l}C_n^2\ \text{种}\\ n-1\ \text{种}\end{array}\right.$ % 原稿左侧“一群”(上)“一个”(下)两行共用一个大括号；右侧大括号上行“$C_n^2$种”、下行“$n-1$种”，分别与“一群”“一个”上下对应

高$\to$低\qquad 低$\to$高（吸收） % 原稿此两处写在上一行下方，“高→低”在左，“低→高(吸收)”偏右

光照跃迁\quad 不多不少\quad $h\nu=|E_{\text{末}}-E_{\text{初}}|$

轰击跃迁\quad $E_k>|E_{\text{末}}-E_{\text{初}}|$（不相等）\quad 只多不少

光照电离\quad 只多不少\quad $h\nu\ge|E_{\text{初}}|=|0-E_{\text{初}}|$（可以相等）

轰击电离\quad 只多不少\quad $E_k>|0-E_{\text{初}}|=|E_{\text{初}}|$
% 原稿左侧有一个大括号把“轰击跃迁、光照电离、轰击电离”三行括在一起（“光照跃迁”在括号之外）
%%%END
%%%BLOCK id=007-02 ch=17 sec=天然放射现象 kind=summary alt=none cont=none y=0.25-0.39
\textbf{原子核衰变过程}\quad 放出$\alpha$、$\beta$、$\gamma$（射线）称天然放射现象（贝克勒尔发现）\quad \textsuperscript{原子}核是复杂的

\begin{center}
\small\setlength{\tabcolsep}{4pt}
\begin{tabular}{c l c c c c c l}
\toprule
 & \textsuperscript{本质} & 质量 & 电量 & 电磁偏转 & $v$ & \makecell{电离性\\ $\alpha$\unclear{m}} & \makecell{穿透性\\ $\alpha$\unclear{$\vee$}} \\
\midrule
$\alpha$ & $\nuc{4}{2}{He}$\ $\alpha$粒子 & $4M$ & $+2e$ & 能 & 小 & \multirow{3}{*}{$\uparrow$\,强} & 纸 \\
$\beta$ & $\nuc{0}{-1}{e}$\ $\beta$粒子 & $\approx0$ & $-e$ & 能 & 中 & & 铝板 \\
$\gamma$ & 光子\ 电磁波 & $0$ & $0$ & 不能 & 大 & & $\downarrow$\,强\ \unclear{水泥墙} \\
\bottomrule
\end{tabular}
\end{center}
% 原稿：“电离性”列下有一条自下而上贯穿三行的长箭头，箭头上端写“强”；“穿透性”列有一条自上而下贯穿三行的长箭头，箭头下端写“强”，三行分别为“纸”“铝板”“水泥墙”（该行末字手写不清）；表格左侧有一条竖线把α、β、γ与右侧隔开
% CHECK: “电离性”“穿透性”两个表头后面的小字（形似“αm”“α∨”）含义不明，照原样标为不清
%%%END
%%%BLOCK id=007-03 ch=17 sec=核反应方程 kind=knowledge alt=none cont=none y=0.385-0.64
质量数$\to\nuc{235}{92}{U}$\ $\leftarrow$电荷量 % 原稿：“质量数”的箭头指向上标235，“电荷量”的箭头由左下方指向下标92

\begin{center}
{\renewcommand{\arraystretch}{1.6}
\begin{tabular}{|c|c|c|}
\hline
$\nuc{1}{1}{H}$ & $\nuc{2}{1}{H}$ & $\nuc{3}{1}{H}$ \\ \hline % 原稿第三格“3H”的下标1不明显
$\nuc{1}{0}{n}$ & $\nuc{0}{-1}{e}$ & $\nuc{4}{2}{He}$ \\ \hline
\end{tabular}}
\end{center}

\textbf{核反应}\quad 三变\ 四恒\ 一\unclear{奇葩}

守恒：质量数、电荷数、能量、动量

不守恒：质子数、中子数、电子数（$\beta$衰变）

质量：质量亏损$\checkmark$\ （静态质量）\quad 质量守恒$\checkmark$\ （相对论质量）\quad $E=mc^2$ % 原稿两处括号写在对应词的下一行

$M_{\text{反}}\neq M_{\text{生}}$
%%%END
%%%BLOCK id=007-04 ch=17 sec=原子核衰变 kind=knowledge alt=11 cont=none y=0.64-0.87
\textbf{衰变}
\begin{itemize}
\item 一变多：天然衰变，自发的
\item 83号以上元素都可以衰变，83号以下元素部分可衰变\quad $\nuc{1}{1}{H}$\quad $\nuc{17}{8}{O}\ \checkmark$\quad $\nuc{16}{8}{O}\ \times$ % “$^{17}_{8}$O”的上标17手写不太清（也可能是18）
\end{itemize}

\textbf{$\alpha$衰变}\quad 本质：核子的重新组合\qquad 都逆时针转

\hspace{2em}一变多\textsuperscript{二}，生成物中有$\alpha$粒子（氦核$\nuc{4}{2}{He}$）\quad 退2少4\ 进1保质\quad $R=\dfrac{p}{qB}$\ 外切 % “多”字上方有一个小“=”（疑为“二”）

\textbf{$\beta$衰变}\quad 本质\ $\nuc{1}{0}{n}\to\nuc{1}{1}{H}+\nuc{0}{-1}{e}$（分裂）\qquad $R\propto\dfrac{1}{q}$ % “$R\propto 1/q$”原稿写在α、β两行右侧之间

\hspace{2em}一变多，生成物中有$\beta$粒子（电子$\nuc{0}{-1}{e}$）\quad 电子顺新核逆时针\quad $R=\dfrac{p}{qB}$\ 内切

\textbf{$\gamma$衰变}\quad $\gamma$射线：电磁波（光子）\textsuperscript{仅能}伴随$\alpha$衰变和$\beta$衰变产生
%%%END
%%%BLOCK id=007-05 ch=17 sec=人工转变 kind=knowledge alt=none cont=none y=0.87-1.00
\textbf{人工转变}
\begin{itemize}
\item 二变二。非自发\textsuperscript{核}反应 % 原稿“二变=.”，末字“二”写作“=”
\item 反应物中有$\alpha$粒子\ $\nuc{4}{2}{He}$
\end{itemize}

\[
\nuc{4}{2}{He}+\nuc{14}{7}{N}\to\nuc{17}{8}{O}+\nuc{1}{1}{H}\qquad\text{卢瑟福发现质子}
\]
\[
\nuc{4}{2}{He}+\nuc{9}{4}{Be}\to\nuc{12}{6}{C}+\nuc{1}{0}{n}\qquad\text{查德威克发现中子}
\]
%%%END
%%%PAGEEND 7
